V tem poglavju bomo najprej predstavili pravila sklepanja, na katerih temelji naravna dedukcija oziroma ``standardna'' logika in s tem bomo tudi predstavili notacijo, ki bo uporabljena v tem diplomskem delu.


\pravilo{Konjunkcija} 
Pravilo vpeljave za konjunkcijo pravi, če sta izjavi \(A\) in \(B\) resnični, potem je tudi izjava \(A\) in \(B\) resnična. To lahko zapišemo kot
%
\begin{prooftree}
    \AXC{\(A\)}
    \AXC{\(B\)}
    \RL{\(\land I\).}
    \BIC{\(A \land B\)}
\end{prooftree}
%
Nad črto zapišemo končno mnogo premis in pod črto zapišemo \emph{sklep}. Desno od črte zapišemo ime pravila, ki ga ponavadi okrajšamo. V našem primeru \(\land\)I predstavlja pravilo vpeljave za konjuncijo (ang.~conjunction introduction). 
Splošna oblika pravila sklepanja je torej
%
\begin{prooftree}
    \AXC{\(\mathcal{H}_1\)}
    \AXC{\(\mathcal{H}_2\)}
    \AXC{\(\ldots\)}
    \AXC{\(\mathcal{H}_n\)}
    \RL{ime.}
    \QuaternaryInfC{\(\mathcal{S}\)}
\end{prooftree}
%
Izjavam \(S, \mathcal{H}_1, \ldots \mathcal{H}_n\) pravimo sodbe~\cite[Section 2]{pfenning2009lecture}.
Sodbam \(\mathcal{H}_1, \ldots \mathcal{H}_n\) pravimo \emph{premise} in  \(\mathcal{S}\) je \emph{sklep}
~\cite{Jazbec2022Primerjava}. 
Kasneje bomo videli, kako lahko taka pravila sestavljamo skupaj, da dobimo izpeljana pravila. Takim zapisom dokazov pravimo \emph{sklepna drevesa}. % TODO: vir
Navadno pišemo zraven še kontekst, ki ga ponavadi označimo z \(\Gamma\). Sodba \(\G \vdash A\) pravi, da iz konteksta \(\G\) lahko vzamemo izjavo \(A\) in jo uporabimo. % \cite homotopy type theory, poglavje 1.1
%
\begin{prooftree}
    \AXC{\(\G \vdash A\)}
    \AXC{\(\G \vdash B\)}
    \RL{\(\land I\)}
    \BIC{\(\G \vdash A \land B\)}
\end{prooftree}
%
Zgornje pravilo preberemo kot, da če imamo poljubni resnični izjavi \(A\) in \(B\) v konekstu \(\G\), potem v tem kontekstu velja tudi \(A \land B\).
Kontekst je pravzaprav nabor izjav \(A_1 , \ldots, A_n\), ki jih imamo za resnične.
Kljub temu, da nabor ponavadi v matematiki pomeni urejeno zaporedje, se tukaj zavedajmo, da zaporedje izjav v kontekstu ni pomembno.

Za konjunkcijo imamo dve pravili uporabe (ang.~elimination rule) --- levo in desno. Levo pravilo iz iz konjunkcije vzame levo stran, torek če imamo \(A \land B\) potem imamo \(A\). Desno pravilo pa vzame desno stran konjunkcije.
%
\begin{equation*}
    \begin{bprooftree}
        \AXC{\(\G \vdash A \land B\)}
        \RL{\(\land E_L\)}
        \UIC{\(\G \vdash A\)}
    \end{bprooftree}
    \qquad
     \begin{bprooftree}
        \AXC{\(\G \vdash A \land B\)}
        \RL{\(\land E_D\)}
        \UIC{\(\G \vdash A\)}
    \end{bprooftree}
\end{equation*}

\pravilo{Implikacija}
Malce bolj zahtevno pravilo vpeljave ima implikacija. Če lahko iz \(A\) izpeljemo \(B\), potem velja \(A \implies B\). To zapišemo kot
%
\begin{prooftree}
    \AXC{\(\G, A \vdash B\)}
    \RL{\(\Rightarrow I\)}
    \UIC{\(\G \vdash A \implies B\)}
\end{prooftree}
%
Torej če lahko iz konteksta \(\G\), ki mu je pridružena izjava \(A\), izpeljemo izjavo \(B\), potem lahko iz konteksta \(\G\) izpeljemo implikacijo \(A \implies B\). 

Da lahko vidiimo pravilo vpeljave na primeru, potrebujemo še emo pravilo sklepanja, in sicer pravilo aksioma
%
\begin{prooftree}
    \AXC{}
    \RL{\emph{Ax}}
    \UIC{\(\G, A \vdash A\)}
\end{prooftree}
%
To pravilo pravi, da brez dodatnih pogojev lahko sklepamo, da iz konteksta, ki vsebuje \(A\) lahko vzamemo \(A\). Bodimo pozorni, da kontekst tu ne pomeni samo \(\Gamma\), ampak \(\Gamma, A\), torej naboru izjav \(\Gamma = (A_1, \ldots, A_n)\) pridružimo še \(A\), da dobimo \(A_1, \ldots A_n, A\). Pogosto se v literaturi pojavlja tudi notacija
%
\begin{prooftree}
    \AXC{\(A \in \G\)}
    \RL{\emph{Ax}.}
    \UIC{\(\G \vdash A\)}
\end{prooftree}

Sedaj si lahko pogledamo primer dokaza za \(A \Rightarrow (B \Rightarrow (A \land B))\)
% ~\cite[stran 6]{pfenning2009lecture}.
%
% TODO: spremeni potek dokaza
\begin{prooftree}
    \AXC{}
    \RL{\emph{Ax}}
    \UIC{\(A, B \vdash A\)}
    \AXC{}
    \RL{\emph{Ax}}
    \UIC{\(A, B \vdash B\)}
    \RL{\(\land I\)}
    \BIC{\(A, B \vdash A \land B\)}
    \RL{\(\Rightarrow I\)}
    \UIC{\(A \vdash B \implies (A \land B)\)}
    \RL{\(\Rightarrow I\)}
    \UIC{\(\vdash A \implies (B \implies (A \land B))\)}
\end{prooftree}
%
Pri tem primeru bi radi poudarili strukturo \emph{dokaznega drevesa}. Pravila dokazovanja se gradijo v večja \emph{izpeljana pravila} v obliki usmerjenega drevesnega grafa, kjer so vozlišča sodbe in povezave pravila med njimi.
Razvidno je tudi, da trditev \(A \Rightarrow (B \Rightarrow (A \land B))\) drži sama po sebi in nismo potrebovali nobenih premis, da smo to pokazali. Pravilo aksioma deluje kot predpostavka. 

% Naj bralca ne zmoti dejstvo, da smo označevali pridružitev izjave h kontekstu ter združitev dveh izjav v kontekst z enako oznako ``\(,\)''. Kontekst namreč predstavlja izjave \(A_1, \ldots A_n\), torej notacija \(\G, A\) predstavlja \(A_1, \ldots A_n, A\). 

Naj bralca ne zmoti pomankanje konteksta v zadnjem sklepu, namreč \(\G\) lahko predstavlja tudi prazen kontekst, torej smo še vedno upoštevali pravila sklepanja.


Pravilo uporabe za implikacijo je enostavnejše.
%
\begin{prooftree}
    \AXC{\(\G \vdash A \implies B\)}
    \AXC{\(\G \vdash A\)}
    \RL{\(\Rightarrow E\)}
    \BIC{\(\G \vdash B\)}
\end{prooftree}
%
Če vemo, da iz \(A\) sledi \(B\), in vemo \(A\), potem vemo \(B\).

\pravilo{Disjunkcija}
Disjunkcija ima levo in desno pravilo vpeljave.
%
\begin{equation*}
    \begin{bprooftree}
        \AXC{\(\G \vdash A\)}
        \RL{\(\lor I_L\)}
        \UIC{\(\G \vdash A \lor B\)}
    \end{bprooftree}
    \begin{bprooftree}
        \AXC{\(\G \vdash B\)}
        \RL{\(\lor I_D\)}
        \UIC{\(\G \vdash A \lor B\)}
    \end{bprooftree}
\end{equation*}
%
Ko v dokazu želimo uporabiti disjunkcijo, ne vemo, ali drži leva ali desna stran. Torej ju ločimo in pokažemo, da iz obeh lahko dobimo željen cilj.
%
\begin{prooftree}
    \AXC{\(\G \vdash A \lor B\)}
%
    \AXC{\(\G, A \vdash C\)}
%
    \AXC{\(\G, B \vdash C\)}
    \RL{\(\lor E\) }
    \TIC{\(\G \vdash C\)}
\end{prooftree}
%
Vidimo da je to pravilo sklepanja podobno pravilu uporabe za implikacijo, saj je spet potrebno pokazazati, da iz \(A\) lahko dobimo \(C\), kot da bi vpeljali implikacijo. 
Poglejmo si primer kako uporabimo disjunkcijo, tako da dokažemo \((A \lor B) \implies (B \lor A)\). Najprej si samo vizualizirajmo, kaj sploh želimo dokazati.
%
\begin{prooftree}
    \AXC{}
    \RL{\emph{Ax}}
    \UIC{\(A \lor B \vdash A \lor B\)}
    \noLine
    \UIC{\(\vdots\)}
    \noLine
    \UIC{\(A \lor B \vdash B \lor A\)}
    \RL{\(\Rightarrow I\)}
    \UIC{\( \vdash (A  \lor B) \implies (B \lor A)\)}
\end{prooftree}
%
Nadaljujemo na edini način,  ki nam je na voljo, in sicer razdelitev \(A \lor B\) na primere.
%
\begin{prooftree}
    \AXC{}
    \RL{\emph{Ax}}
    \UIC{\(A \lor B \vdash A \lor B\)}
    \AXC{?}
    \UIC{\(A \lor B, A \vdash B \lor A\)}
    \AXC{?}
    \UIC{\(A \lor B, B \vdash B \lor A\)}
    \RL{\(\lor E\)}
    \TIC{\(A \lor B \vdash B \lor A\)}
    \RL{\(\Rightarrow I\)}
    \UIC{\( \vdash (A  \lor B) \implies (B \lor A)\)}
\end{prooftree}
%
Moramo samo še ugotoviti, kaj napisati namesto vprašajev.
Lotimo se manjšega podprimera.
\begin{prooftree}
    \AXC{?}
    \UIC{\(A \lor B, A \vdash B \lor A\)}
\end{prooftree}
%
Glede na to, da v sklepu potrebujemo dijunkcijo, lahko poskusimo vpeljati to disjunkcijo. 
\begin{prooftree}
    \AXC{?}
    \UIC{\(A \lor B, A \vdash A\)}
    \RL{\(\lor I_D\)}
    \UIC{\(A \lor B, A \vdash B \lor A\)}
\end{prooftree}
%
Da iz konteksta \(A \lor B, A\) lahko vzamemo \(A\) pa nam pove pravilo hipoteze.
%
\begin{prooftree}
    \AXC{}
    \RL{\emph{Ax}}
    \UIC{\(A \lor B, A \vdash A\)}
    \RL{\(\lor I_D\)}
    \UIC{\(A \lor B, A \vdash B \lor A\)}
\end{prooftree}
%
Simetrično dobimo dokaz 
\begin{prooftree}
    \AXC{}
    \UIC{\(A \lor B, B \vdash B \lor A\)}
\end{prooftree}
%
Vstavimo to v naš večji dokaz, da dobimo celotno izpeljavo z osnovnimi pravili.
%
\begin{prooftree}
    \AXC{}
    \RL{\emph{Ax}}
    \UIC{\(A \lor B \vdash A \lor B\)}

    \AXC{}
    \RL{\emph{Ax}}
    \UIC{\(A \lor B, A \vdash A\)}
    \RL{\(\lor I_D\)}
    \UIC{\(A \lor B, A \vdash B \lor A\)}

    \AXC{}
    \RL{\emph{Ax}}
    \UIC{\(A \lor B, B \vdash B\)}
    \RL{\(\lor I_L\)}
    \UIC{\(A \lor B, B \vdash B \lor A\)}

    \RL{\(\lor E\)}
    \TIC{\(A \lor B \vdash B \lor A\)}
    \RL{\(\Rightarrow I\)}
    \UIC{\( \vdash (A  \lor B) \implies (B \lor A)\)}
\end{prooftree}

V kasnejših poglavjih bomo povedali še pravila za eksistenčni in univerzalni kvantifikator ter za negacijo. 


% \pravilo{Negacija}
% Za negacijo imamo več možnosti osnovnih pravil. Navadno se negacija \(\neg A\) predstavi kot \(A \implies \bot\), se pravi, če iz \(A\) lahko pokažemo neresnico, potem \(A\) ne velja. To podre strukturo pravil vpeljave in uporabe

% \pravilo{Resnica in neresnica}


% pravila za konjunkcijo, impolikacijo, disjunkcijo (cases), negacija, resnica
% univerzalni kvantifikator, [t/x]A
% eliminacija, vpeljava in klasična logika
% kontekst, konkatenacija, razlaga notacije
